\documentclass[a4paper,11pt,reqno]{amsart}
\usepackage[margin=1.2in,marginparwidth=1.5cm,marginparsep=0.5cm]{geometry}
\usepackage{graphicx}
\usepackage{amsmath,amssymb,amsthm}
\usepackage{mathtools}

\title{Exercícios de ``Introdução às Equações Diferenciais''  IMPA-TECH 2026-01}

\begin{document}

\maketitle

\section*{Banco de Questões}

\begin{enumerate}

\item \textbf{Problemas de valor inicial.}
Resolva os seguintes problemas de valor inicial. Em cada item, indique explicitamente o maior intervalo em que a solução está definida.

\begin{enumerate}
    \item 
    \[
        y' + 2y = e^{-t}, \qquad y(0)=1.
    \]

    \item 
    \[
        y' = t(1+y^2), \qquad y(0)=0.
    \]

    \item 
    \[
        y' = \frac{t+y}{t}, \qquad y(1)=0, \qquad t>0.
    \]
    \emph{Dica:} tente reconhecer uma expressão simples cuja derivada aparece na equação.
\end{enumerate}


\item \textbf{Teoria qualitativa e unicidade.}
Considere a equação autônoma
\[
    x' = x(x-1)(x-2).
\]

\begin{enumerate}
    \item Encontre todas as soluções de equilíbrio.

    \item Determine o sinal de $x'$ em cada um dos intervalos
    \[
        (-\infty,0), \qquad (0,1), \qquad (1,2), \qquad (2,\infty),
    \]
    e desenhe a reta de fase.

    \item Suponha que $x(0)=x_0$. Sem resolver explicitamente a equação, determine
    \[
        \lim_{t\to +\infty} x(t)
    \]
    em cada um dos casos
    \[
        x_0<0, \qquad 0<x_0<1, \qquad 1<x_0<2, \qquad x_0>2.
    \]

    \item Prove que, se $x_0\in(0,2)$, então a solução está definida para todo $t\geq 0$.
\end{enumerate}


\item \textbf{Redução de ordem a partir de uma solução conhecida.}
Considere a equação diferencial
\[
    t^2 y'' - 2t y' + 2y = t^3, \qquad t>0.
\]

\begin{enumerate}
    \item Verifique que $y_1(t)=t$ é uma solução da equação homogênea associada.

    \item Encontre uma segunda solução $y_2$ da equação homogênea que seja linearmente independente de $y_1$.

    \item Encontre a solução geral da equação não homogênea.

    \item Resolva o problema de valor inicial
    \[
        t^2 y'' - 2t y' + 2y = t^3, \qquad y(1)=0, \qquad y'(1)=1.
    \]
\end{enumerate}


\item \textbf{Zeros de soluções e o Wronskiano.}
Sejam $p,q:I\to\mathbb{R}$ funções contínuas em um intervalo $I$, e considere a equação linear
\[
    y'' + p(t)y' + q(t)y = 0.
\]

\begin{enumerate}
    \item Se $y_1,y_2$ são duas soluções, prove a identidade de Abel:
    \[
        W(t)=W(t_0)\exp\left(-\int_{t_0}^t p(s)\,ds\right),
    \]
    onde
    \[
        W(t)=
        \begin{vmatrix}
            y_1(t) & y_2(t)\\
            y_1'(t) & y_2'(t)
        \end{vmatrix}.
    \]

    \item Suponha que $y$ seja uma solução não identicamente nula. Prove que $y$ não pode ter um zero duplo; isto é, não existe $t_0\in I$ tal que
    \[
        y(t_0)=0
        \qquad\text{e}\qquad
        y'(t_0)=0.
    \]

    \item Suponha que $y_1,y_2$ sejam soluções linearmente independentes e que
    \[
        y_1(t_0)=y_2(t_0)=0
    \]
    para algum $t_0\in I$. Mostre que isso é impossível.

    \item Dê um exemplo de uma equação linear de segunda ordem e de duas soluções linearmente independentes $y_1,y_2$ tais que $y_1$ e $y_2$ tenham zeros, mas nunca no mesmo ponto.
\end{enumerate}


\item \textbf{Reconstrução de uma equação a partir do espaço de soluções.}
Seja $I\subset\mathbb{R}$ um intervalo, e sejam $u,v\in C^2(I)$ duas funções tais que
\[
    W(u,v)(t)=
    \begin{vmatrix}
        u(t) & v(t)\\
        u'(t) & v'(t)
    \end{vmatrix}
    \neq 0
    \qquad\text{para todo }t\in I.
\]

\begin{enumerate}
    \item Prove que existem únicas funções contínuas $p,q:I\to\mathbb{R}$ tais que $u$ e $v$ são ambas soluções de
    \[
        y'' + p(t)y' + q(t)y = 0.
    \]

    \item Encontre fórmulas explícitas para $p(t)$ e $q(t)$ em termos de $u,v$ e suas derivadas.

    \item Aplique o resultado a
    \[
        u(t)=e^t,
        \qquad
        v(t)=te^t,
        \qquad
        I=\mathbb{R}.
    \]
    Encontre a única equação da forma
    \[
        y''+p(t)y'+q(t)y=0
    \]
    para a qual $\{e^t,te^t\}$ é uma base do espaço de soluções.

    \item Agora suponha apenas que $u,v\in C^2(I)$ sejam linearmente independentes, mas que $W(u,v)$ se anule em algum ponto de $I$. Explique por que, nesse caso, não podem existir funções contínuas $p,q:I\to\mathbb{R}$ tais que $u$ e $v$ formem uma base do espaço de soluções de uma equação
    \[
        y''+p(t)y'+q(t)y=0
    \]
    em todo o intervalo $I$.
\end{enumerate}


\item \textbf{Lei de resfriamento.}
A taxa de variação da temperatura $T(t)$ de um corpo é proporcional à diferença entre a temperatura do corpo e a temperatura do ambiente $T_m$, em que $k > 0$ denota a constante de proporcionalidade.

Um componente metálico de um motor é retirado de um forno a uma temperatura de $100^\circ C$ e colocado em um galpão onde a temperatura ambiente é mantida constante a $20^\circ C$. Após $10$ minutos, a temperatura do componente caiu para $60^\circ C$.

\begin{enumerate}
    \item Determine a expressão da temperatura $T(t)$ em qualquer instante $t$.
    \item Quanto tempo, aproximadamente, levará, desde o instante em que foi retirado do forno, para que a temperatura do componente atinja $30^\circ C$?
\end{enumerate}


\item \textbf{Equação de Legendre.}
Determine duas soluções linearmente independentes de
\[
    (1 - t^2)\ddot{x} - 2t\dot{x} + \lambda(\lambda + 1)x = 0,
    \qquad -1 < t < 1.
\]

\emph{Sugestão:} use séries de potências centradas em $t = 0$. Mostre que, quando $\lambda$ é inteiro $\geq 0$, uma dessas soluções é um polinômio.


\item \textbf{Oscilador com forçamento especial.}
Considere um sistema físico cuja posição $y(t)$ é regida pela EDO de segunda ordem
\[
    y'' + 4y = \sec(2t)
\]
com as condições iniciais de repouso
\[
    y(0)=0,
    \qquad
    y'(0)=0.
\]

\begin{enumerate}
    \item Determine a solução geral da equação homogênea associada.
    \item Utilizando o método de variação de parâmetros, determine uma solução particular $y_p(t)$ para o intervalo $0\leq t<\pi/4$.
    \item Resolva o problema de valor inicial e expresse $y(t)$ de forma simplificada.
    \item Discuta o que acontece com o deslocamento $y(t)$ à medida que $t$ se aproxima de $\pi/4$.
\end{enumerate}


\item \textbf{Existência e unicidade.}
Considere o problema de valor inicial
\[
    \frac{dy}{dt}=3y^{2/3},
    \qquad
    y(t_0)=y_0.
\]

\begin{enumerate}
    \item No ponto $(t_0,y_0)=(0,1)$, determine se existe solução e se ela é única. Caso não seja, exiba duas soluções.
    \item Faça o mesmo para o ponto $(t_0,y_0)=(0,0)$.
    \item Explique por que os dois casos acima têm comportamentos diferentes à luz do teorema de existência e unicidade.
\end{enumerate}


\item \textbf{Equação linear de ordem $n$.}
Sejam $a_0,\ldots,a_{n-1}$ funções contínuas reais, ou complexas, em um intervalo $I$.

A equação linear
\[
\frac{d^n x}{dt^n}
=
a_{n-1}(t)\frac{d^{n-1}x}{dt^{n-1}}
+\cdots+a_0(t)x
\tag{*}
\]
é chamada de equação linear de ordem $n$.

Considere $C^n=C^n(I,\mathbb{R})$, ou $C^n(I,\mathbb{C})$, o espaço vetorial das funções reais, ou complexas, de classe $C^n$ em $I$.

Prove que:

\begin{enumerate}
    \item O conjunto das soluções de $(*)$ é um subespaço vetorial de $C^n$ de dimensão $n$.

    \emph{Sugestão:} escreva $x_1=x$, $x_2=x'$, \ldots, $x_n=x^{(n-1)}$, e verifique que $(*)$ é equivalente a um sistema linear da forma
    \[
        X'=A(t)X.
    \]

    \item Sejam $\varphi_1,\ldots,\varphi_n\in C^n$ soluções de $(*)$ e seja
    \[
        W(t)=W(\varphi_1,\ldots,\varphi_n)(t)
    \]
    o determinante da matriz $n\times n$ cuja $i$-ésima linha é formada por
    \[
        \frac{d^{i-1}\varphi_1}{dt^{i-1}},
        \ldots,
        \frac{d^{i-1}\varphi_n}{dt^{i-1}},
    \]
    para $i=1,\ldots,n$. Prove que
    \[
        W(t)=W(t_0)\exp\left(\int_{t_0}^t a_{n-1}(s)\,ds\right).
    \]
    Conclua que $n$ soluções $\varphi_1,\ldots,\varphi_n$ são linearmente independentes se, e somente se,
    \[
        W(\varphi_1,\ldots,\varphi_n)(t)\neq 0
        \qquad\text{para todo }t\in I.
    \]

    \item Sejam $\varphi_1,\ldots,\varphi_n\in C^n$ tais que
    \[
        W(\varphi_1,\ldots,\varphi_n)(t)\neq 0
        \qquad\text{em }I.
    \]
    Prove que existe uma única equação da forma $(*)$ que tem $\{\varphi_1,\ldots,\varphi_n\}$ como base do espaço de soluções.
\end{enumerate}


\item \textbf{Existência de uma função.}
Existe uma função não nula $f:\mathbb{R}\to\mathbb{R}$ tal que
\[
    f^{(4)}(t)+f(t)=0
\]
para todo $t\in\mathbb{R}$, e
\[
    \lim_{t\to+\infty}f(t)
    =
    \lim_{t\to+\infty}f'(t)
    =
    0?
\]


\item \textbf{Equação com soluções prescritas.}
Suponha que as funções $\sin t$ e $\sin 2t$ sejam ambas soluções de uma equação diferencial linear homogênea com coeficientes constantes
\[
    x^{(n)}+c_{n-1}x^{(n-1)}+\cdots+c_1x'+c_0x=0,
\]
com
\[
    c_0,c_1,\ldots,c_{n-1}\in\mathbb{R}.
\]
Qual é a menor ordem que essa equação pode ter? Escreva uma equação diferencial de ordem mínima que tenha ambas as funções como soluções.


\item \textbf{Condição necessária para existência de solução.}
Sejam $a:[0,1]\to(0,+\infty)$ uma função contínua e $\lambda\in\mathbb{R}$. Mostre que, se a equação diferencial
\[
    x''+\lambda a(t)x=0
\]
admite uma solução não identicamente nula satisfazendo
\[
    x(0)=0,
    \qquad
    x'(1)=0,
\]
então $\lambda>0$.

Exiba um exemplo com $\lambda>0$ para o qual existe uma solução sob essas condições.


\item \textbf{Soluções maximais escapam de compactos.}
Seja $f:U\subset\mathbb{R}^{n+1}\to\mathbb{R}^n$ uma função de classe $C^1$, definida no aberto $U$, e seja $(t_0,x_0)\in U$.
Suponha que $x:(\alpha,\beta)\to\mathbb{R}^n$ seja uma solução maximal de
\[
    x'(t)=f(t,x(t)),
    \qquad
    x(t_0)=x_0,
\]
em $U$.
Se $\beta<\infty$, mostre que, dado qualquer compacto $K\subset U$, existe $t^*$ com
\[
    t_0<t^*<\beta
\]
tal que
\[
    (t^*,x(t^*))\in U\setminus K.
\]
O enunciado análogo vale se $-\infty<\alpha$.


\item \textbf{Ondas viajantes para a equação do calor linear.}
Considere a equação do calor linear
\begin{equation}\label{linear-heat-equation}
    \partial_t u-\partial_x^2u=0,
\end{equation}
onde $x,t\in\mathbb{R}$ e uma solução é uma função $u\in C^2(\mathbb{R}^2)$ satisfazendo a equação.

Dizemos que uma solução $u\in C^2(\mathbb{R}^2)$ é uma onda viajante para a equação do calor se existe uma função $\phi\in C^2(\mathbb{R})$ tal que
\[
    u(x,t)=\phi(x-ct)
\]
para algum $c\in\mathbb{R}$. Se, além disso, $\phi$ satisfaz
\[
    \phi(x)\to 0
    \quad\text{quando }x\to-\infty,
    \qquad
    \phi(x)\to x_\infty
    \quad\text{quando }x\to+\infty
\]
para algum $x_\infty\in\mathbb{R}$, dizemos que $u$ é um kink.

\begin{enumerate}
    \item Se $u(x,t)=\phi(x-ct)$ é uma onda viajante para a equação do calor linear, que equação a função $\phi$ deve satisfazer?

    \emph{Dica:} substitua $u$ na equação \eqref{linear-heat-equation} e encontre uma equação diferencial para $\phi$.

    \item Se $u(x,t)=\phi(x-ct)$ é um kink para a equação do calor linear, que outra equação a função $\phi$ deve satisfazer?

    \emph{Dica:} integre a equação obtida no item anterior de $-\infty$ até $x$.

    \item Existem kinks para a equação do calor linear? É possível escrever a fórmula de um kink?
\end{enumerate}


\item \textbf{Ondas viajantes para a equação do calor não linear.}
Considere a equação do calor não linear
\begin{equation}\label{nonlinear-heat-equation}
    \partial_t u-\partial_x^2u=\frac12\partial_x(u^2),
\end{equation}
onde $x,t\in\mathbb{R}$ e uma solução é uma função $u\in C^2(\mathbb{R}^2)$ satisfazendo a equação.

\begin{enumerate}
    \item Se $u(x,t)=\phi(x-ct)$ é uma onda viajante para a equação do calor não linear, que equação a função $\phi$ deve satisfazer?

    \item Se $u(x,t)=\phi(x-ct)$ é um kink para a equação do calor não linear, que outra equação a função $\phi$ deve satisfazer?

    \item Existem kinks para a equação do calor não linear? É possível escrever a fórmula de um kink?
\end{enumerate}


\item \textbf{Problemas de valor inicial.}
Resolva os seguintes problemas de valor inicial. Em cada item, indique explicitamente o maior intervalo em que a solução está definida.

\begin{enumerate}
    \item
    \[
        y'-3y=6e^{3t},
        \qquad
        y(0)=1.
    \]

    \item
    \[
        y'=\frac{t}{y},
        \qquad
        y(0)=2.
    \]

    \item
    \[
        y'=\frac{2t+y}{t},
        \qquad
        y(1)=1,
        \qquad
        t>0.
    \]
    \emph{Dica:} tente escrever a equação em termos de $y/t$.
\end{enumerate}


\item \textbf{Equação exata com condição inicial.}
Considere a equação diferencial
\[
    (2xy+y^2)\,dx+(x^2+2xy)\,dy=0.
\]

\begin{enumerate}
    \item Verifique que a equação é exata.

    \item Encontre uma função $F(x,y)$ tal que
    \[
        dF=(2xy+y^2)\,dx+(x^2+2xy)\,dy.
    \]

    \item Encontre a solução implícita do problema de valor inicial
    \[
        (2xy+y^2)\,dx+(x^2+2xy)\,dy=0,
        \qquad
        y(1)=1.
    \]

    \item Determine se é possível escrever a solução explicitamente como $y=y(x)$ em uma vizinhança de $x=1$.
\end{enumerate}


\item \textbf{Existência, unicidade e soluções maximais.}
Considere o problema de valor inicial
\[
    y'=y^2,
    \qquad
    y(0)=y_0.
\]

\begin{enumerate}
    \item Resolva explicitamente o problema de valor inicial para $y_0\neq 0$.

    \item Determine o intervalo maximal de existência da solução em função de $y_0$.

    \item Explique por que não há contradição entre o fato de a função $f(t,y)=y^2$ ser suave em todo $\mathbb{R}^2$ e o fato de algumas soluções não estarem definidas para todo $t\in\mathbb{R}$.

    \item Descreva o comportamento da solução quando ela se aproxima de uma extremidade finita do seu intervalo maximal.
\end{enumerate}


\item \textbf{Dependência linear e zeros do Wronskiano.}
Considere duas funções $u,v\in C^2(I)$ em um intervalo $I\subset\mathbb{R}$.

\begin{enumerate}
    \item Mostre que, se $u$ e $v$ são soluções de uma mesma equação linear homogênea
    \[
        y''+p(t)y'+q(t)y=0,
    \]
    com $p,q$ contínuas em $I$, então ou
    \[
        W(u,v)(t)=0
        \qquad\text{para todo }t\in I,
    \]
    ou
    \[
        W(u,v)(t)\neq 0
        \qquad\text{para todo }t\in I.
    \]

    \item Dê um exemplo de duas funções $u,v\in C^2(\mathbb{R})$ linearmente independentes tais que
    \[
        W(u,v)(0)=0.
    \]

    \item Explique por que as funções do item anterior não podem ser duas soluções linearmente independentes de uma mesma equação
    \[
        y''+p(t)y'+q(t)y=0
    \]
    em todo $\mathbb{R}$, com $p,q$ contínuas.

    \item Discuta o que isso mostra sobre a diferença entre ``funções linearmente independentes'' e ``soluções linearmente independentes de uma mesma EDO linear homogênea''.
\end{enumerate}


\item \textbf{Problema de contorno impossível.}
Seja $q:[0,1]\to\mathbb{R}$ uma função contínua tal que
\[
    q(t)\leq 0
    \qquad\text{para todo }t\in[0,1].
\]
Considere a equação
\[
    y''+q(t)y=0.
\]

\begin{enumerate}
    \item Suponha que $y$ seja uma solução tal que
    \[
        y(0)=0
        \qquad\text{e}\qquad
        y(1)=0.
    \]
    Multiplique a equação por $y$ e integre de $0$ a $1$. Mostre que
    \[
        \int_0^1 (y'(t))^2\,dt
        =
        \int_0^1 q(t)y(t)^2\,dt.
    \]

    \item Conclua que a única solução satisfazendo
    \[
        y(0)=y(1)=0
    \]
    é a solução identicamente nula.

    \item Agora suponha que $q(t)<0$ para todo $t\in[0,1]$. Mostre que não existe solução não identicamente nula satisfazendo
    \[
        y'(0)=0
        \qquad\text{e}\qquad
        y'(1)=0.
    \]

    \item Explique por que esse resultado é surpreendente do ponto de vista de problemas de valor inicial, mas natural do ponto de vista de problemas de contorno.
\end{enumerate}


\item \textbf{Wronskiano, domínio e reconstrução de uma equação.}
Seja $I\subset\mathbb{R}$ um intervalo aberto e seja
\[
    L:C^2(I)\to C^0(I)
\]
um operador da forma
\[
    L[y](t)=y''(t)+p(t)y'(t)+q(t)y(t),
\]
onde $p,q:I\to\mathbb{R}$ são funções contínuas.

Mostre que não existe uma equação linear homogênea de segunda ordem com coeficientes contínuos em todo $\mathbb{R}$,
\[
    y''+p(t)y'+q(t)y=0,
\]
que possua simultaneamente as duas funções
\[
    y_1(t)=\sin(t^2),
    \qquad
    y_2(t)=\cos(t^2)
\]
como soluções em todo $\mathbb{R}$.

Mais precisamente:

\begin{enumerate}
    \item Calcule o Wronskiano
    \[
        W(y_1,y_2)(t).
    \]

    \item Explique por que $y_1$ e $y_2$ são linearmente independentes em $\mathbb{R}$.

    \item Suponha, por absurdo, que existam funções contínuas $p,q:\mathbb{R}\to\mathbb{R}$ tais que $y_1$ e $y_2$ sejam soluções de
    \[
        y''+p(t)y'+q(t)y=0
    \]
    em todo $\mathbb{R}$. Use a identidade de Abel para obter uma contradição.

    \item Agora remova o ponto problemático: mostre que, em cada um dos intervalos
    \[
        (-\infty,0)
        \qquad\text{e}\qquad
        (0,+\infty),
    \]
    existe uma única equação da forma
    \[
        y''+p(t)y'+q(t)y=0
    \]
    para a qual $\sin(t^2)$ e $\cos(t^2)$ formam uma base de soluções.

    \item Encontre explicitamente essa equação em $(0,+\infty)$ e em $(-\infty,0)$.
\end{enumerate}

\end{enumerate}

\end{document}